

We provide fully symbolical expressions of the amplitude for double Higgs boson production in gluon fusion mediated by light quarks~\cite{Bonetti:2025vfd}.


\begin{figure}
\centering
    \subfloat[]{{\includegraphics[height=0.11\textheight]{Bonetti_Rendler_Torres_Bobadilla/Images/gg-HHH}}\label{fig:VVHV}}
	\qquad\qquad
    \subfloat[]{{\includegraphics[height=0.11\textheight]{Bonetti_Rendler_Torres_Bobadilla/Images/gg-HHVV.pdf}}\label{fig:HHVV}}
	\qquad\qquad
    \subfloat[]{{\includegraphics[height=0.11\textheight]{Bonetti_Rendler_Torres_Bobadilla/Images/gg-VVV}}\label{fig:VVV}}
	\caption{Representatives for the light-quark contributions.}
    \label{fig:amplgg}
\end{figure}


This mode of production appears for the first time at two loops: the gluons couple to a light-quark box, to which either two $W$ or two $Z$ bosons are attached, subsequently generating the Higgs boson(s). We divide the relevant diagrams into three blocks: one-particle reducible diagrams containing one vector-vector-Higgs and one triple-Higgs vertex (dubbed $VVh$, see Fig.~\ref{fig:VVHV}); diagrams containing a vector-vector-Higgs-Higgs vertex ($VVhh$, Fig.~\ref{fig:HHVV}); and diagrams containing two vector-vector-Higgs vertices ($VVV$, Fig.~\ref{fig:VVV}). Each one of these blocks is further divided into diagrams containing either $W$ or $Z$ bosons only. Each one of the aforementioned terms is explicitly gauge-invariant and UV- and IR-finite.

No axial terms are generated by the fermion loop thanks to charge-parity conservation. The amplitude can then be decomposed as a linear combination of the standard tensors $T_1^{\mu\nu}$ and $T_2^{\mu\nu}$ already used for double Higgs production when considering QCD corrections~\cite{Borowka:2016ehy}. We consider the first two generations of light quarks coupling to the $W$ boson and five light quarks coupling to the $Z$ boson (all but the top quark).

We extract the two form factors and reduce the set of scalar two-loop Feynman integrals appearing in the amplitude to master integrals. We derive differential equations for specific linear combinations of the master integrals w.r.t.\ the Mandelstam variables and the masses, such that their solution can be written in terms of iterated integrals over logarithmic kernels. We improve on this choice of basis by applying a series of rotations to eliminate redundant integration kernels. We fix the integration constants by matching the expressions of the master integrals to their large-$m_{W,Z}$ expansion.

We employ the results to obtain analytic expressions for the form factors, and we further identify independent functions therein for which we build dedicated differential equations for optimized numerical evaluation, which we solve in terms of generalized series expansions, evolving from the large-mass limit to arbitrary points in the physical region of the phase space. We provide a code implementation for the numerical evaluation of our functions based on the \texttt{Mathematica} package \texttt{DiffExp}~\cite{Hidding:2020ytt}. Grids for the light-quark amplitude, as well as for their interference with the LO amplitude have been implemented in the POWHEG-BOX framework or are available upon request.

We observe that the interference between the LO and the light-quarks first form factor represents the main source of contributions except for the high-energy tail, where the second form factor becomes relevant, although very small overall. We notice large cancellations when considering the interference between LO triangle and box diagrams with each of the $VVH$, $VVHH$, and $VVV$ blocks of diagrams, suggesting that the light-quark terms are highly sensitive to variations of the tri-linear Higgs coupling.

The results have been checked against the independent evaluation~\cite{Bhattacharya_TBA}.

~

We present the impact of Yukawa~\cite{Heinrich:2024dnz} and light-quark corrections~\cite{BoHeReToBa_TBA} to the Higgs-pair invariant mass distribution and the transverse momentum distribution in Fig.~\ref{fig:mHH} and \ref{fig:pTH}, respectively.


\begin{figure}
\centering
    \subfloat[]{{\includegraphics[width=0.48\textwidth]{Bonetti_Rendler_Torres_Bobadilla/Images/16_mhh2_Yukawa_s136}}\label{fig:mHH}}
	\quad
    \subfloat[]{{\includegraphics[width=0.48\textwidth]{Bonetti_Rendler_Torres_Bobadilla/Images/17_pt_Yukawa_s136}}\label{fig:pTH}}
	\caption{Light-quark and Yukawa contributions to the Higgs pair invariant mass and Higgs transverse momentum distributions.}
    \label{fig:obs}
\end{figure}

In the Higgs-pair invariant mass distribution both classes of contributions affect mainly the low-energy region, where the Yukawa terms produce an enhancement of almost $+40\%$, while the light-quarks have a negative effect of more than $-15\%$. Both contributions vanish at high energies. While the light quark have a constant negative effect on the distribution, the Yukawa contributions become slightly negative only between $400~\textup{GeV}$ and $800~\textup{GeV}$. Their combined effect produces an increase of more than $+20\%$ at the production threshold, rapidly approaching zero as the energy grows.

A similar behaviour for light quarks is observed in the transverse momentum distribution, where they have a negative effect at low energies of around $-4\%$, slowly approaching zero for increasing energies. The Yukawa contributions have an oscillatory pattern: they raise to $+5\%$ between $50\,\textup{GeV}$ and $150\,\textup{GeV}$, become negative between $200\,\textup{GeV}$ and $400\,\textup{GeV}$, and slowly rise again at higher energies.


